% ============================================================
%  Preamble format Skripsi Politeknik Statistika STIS
%  Disesuaikan dengan Pedoman Skripsi KS Edisi Keenam (2025)
%  Engine: XeLaTeX (fontspec)
% ============================================================

% --- Font: Times New Roman 12 pt ---
%     Pakai berkas Times New Roman ASLI bila tersedia di sistem (mis. Windows,
%     macOS, atau Linux dengan paket msttcorefonts). Bila tidak ada, jatuh ke
%     TeX Gyre Termes (klon metrik Times yang identik dimensinya) agar dokumen
%     tetap dapat dikompilasi di lingkungan mana pun (mis. TinyTeX bawaan).
\usepackage{fontspec}
\IfFontExistsTF{Times New Roman}
  {\setmainfont{Times New Roman}}
  {\setmainfont{TeX Gyre Termes}}

% --- Margin: kiri/dalam 4 cm; atas, kanan/luar, bawah 3 cm.
%     Karena dokumen dicetak bolak-balik (twoside), margin "kiri 4 cm"
%     berarti margin DALAM (binding) 4 cm dan LUAR 3 cm; pada halaman
%     genap posisi otomatis dicerminkan. ---
\usepackage[a4paper,inner=4cm,top=3cm,outer=3cm,bottom=3cm]{geometry}

% --- Spasi: default dokumen 1,5 (dipakai halaman muka, daftar, dan daftar
%     pustaka). TEKS ISI (Bab I dst.) diset 2 spasi sesuai pedoman hlm. 34
%     melalui \doublespacing di awal bagian isi (lihat index.qmd). ---
\usepackage{setspace}
\onehalfspacing
% Unit "1 spasi" (tinggi baris TUNGGAL pada 12 pt) untuk jarak presisi di halaman
% muka sesuai pedoman (mis. 5 spasi, 4,5 spasi, 1,5 spasi). Selaras dgn Overleaf.
\newlength{\spasi}
\AtBeginDocument{\begingroup\setstretch{1}\global\spasi=\baselineskip\endgroup}
% Cegah perataan vertikal (flushbottom) meregang glue antar-elemen (teorema,
% bukti, tabel, gambar) saat halaman tak penuh; pakai bawah ragged (konsisten).
\raggedbottom
% Kotak bantu untuk menyamakan lebar blok tabel/gambar (agar Sumber & judul
% rata dengan tepi kiri tabel/gambar; blok tetap di tengah halaman).
\newsavebox{\skripsiblokbox}

% --- Jarak persamaan (display math) sesuai pedoman hlm. 8 ---
% Pedoman meminta persamaan pada BARIS TERSENDIRI dengan jarak ~1 baris (sama
% dengan jarak antar baris) dari teks sebelum & sesudahnya. Agar tercapai:
%   (1) tulis $$...$$ MENGALIR di dalam paragraf (tanpa baris kosong) - lihat
%       contoh di bab Metode; bila diberi baris kosong, persamaan menjadi
%       paragraf tersendiri dan jaraknya melebar (glue antar-paragraf spasi
%       ganda yang tidak dipengaruhi display-skip), dan
%   (2) skip display diset kecil & seragam di awal SETIAP lingkungan persamaan
%       (di bawah, lewat \AtBeginEnvironment) sehingga jarak rapat = jarak baris.
\usepackage{etoolbox}
\makeatletter
\newcommand{\skripsi@dispskip}{%
  \setlength{\abovedisplayskip}{4\p@ \@plus2\p@ \@minus1\p@}%
  \setlength{\belowdisplayskip}{4\p@ \@plus2\p@ \@minus1\p@}%
  \setlength{\abovedisplayshortskip}{4\p@ \@plus2\p@ \@minus1\p@}%
  \setlength{\belowdisplayshortskip}{4\p@ \@plus2\p@ \@minus1\p@}%
}
\AtBeginEnvironment{equation}{\skripsi@dispskip}
\AtBeginEnvironment{equation*}{\skripsi@dispskip}
\AtBeginEnvironment{align}{\skripsi@dispskip}
\AtBeginEnvironment{align*}{\skripsi@dispskip}
\AtBeginEnvironment{gather}{\skripsi@dispskip}
\AtBeginEnvironment{gather*}{\skripsi@dispskip}
\AtBeginEnvironment{multline}{\skripsi@dispskip}
\makeatother

% --- Tautan dicetak hitam tanpa kotak (sesuai gaya skripsi cetak) ---
% Tautan tetap dapat diklik di PDF, namun tidak berwarna agar konsisten dengan
% contoh pedoman (teks hitam). Mahasiswa dapat mengubahnya di preamble-tambahan.
\AtBeginDocument{\hypersetup{hidelinks}}

% --- Indentasi paragraf & lain-lain ---
\usepackage{indentfirst}
\setlength{\parindent}{1.27cm} % indentasi 1,27 cm (0,5 inci)
% Quarto memuat paket parskip (parskip=6pt, parindent=0). Kembalikan ke gaya
% pedoman: alinea HANYA ditandai indentasi 1 tab, TANPA spasi ekstra antar
% paragraf, sehingga jarak antar paragraf = jarak antar baris (2 spasi).
\setlength{\parskip}{0pt plus 0.5pt}
\setlength{\mathindent}{1.27cm} % persamaan (fleqn) menjorok 1 tab dari batas kiri

% --- Dukungan tabel sesuai pedoman ---
\usepackage{float}        % opsi penempatan [H]
\usepackage{booktabs}     % garis tabel rapi (toprule/midrule/bottomrule)
\usepackage{longtable}    % tabel lintas-halaman ("Tabel X (lanjutan)")
\usepackage{array}
\usepackage{tabularx}     % kolom 'X' yang otomatis pas ke \textwidth & membungkus teks (tabel lebar)
\usepackage{xltabular}    % longtable + kolom 'X' (opsi lebar untuk tabel panjang)
\usepackage{chngcntr}     % \counterwithout: nomor tabel/gambar/persamaan deret tunggal (tak reset per bab)
\usepackage{ragged2e}
\usepackage{enumitem}      % daftar bernomor/berbutir yang dapat disesuaikan (a., 1.)
\usepackage{pdflscape}    % memutar HALAMAN PDF (tabel/gambar lebar dibaca di layar)
\usepackage{rotating}     % memutar ISI gambar/tabel 90° (sidewaysfigure/sidewaystable);
                          % pada dokumen twoside, arah putar OTOMATIS menyesuaikan
                          % halaman ganjil/genap sehingga kepala gambar menghadap tepi luar.

% --- Pengaman agar isi tidak menembus margin kiri/kanan (lebar teks = 14 cm) ---
\usepackage{fvextra}      % melipat baris kode panjang (dipakai bersama 'code-overflow: wrap')
\fvset{breaklines=true,breakanywhere=true}  % lipat baris kode, termasuk di tengah token panjang (URL/path)
\setlength{\emergencystretch}{3em}          % beri ruang regang agar paragraf tak meluber ke kanan
\usepackage{lastpage}     % \pageref{LastPage}: jumlah halaman terakhir (untuk hitung halaman di Abstrak)
\usepackage{refcount}     % \getpagerefnumber{...}: ambil NILAI nomor halaman (untuk dihitung)
% Catatan:
% - Gambar: Quarto membungkus dengan \pandocbounded sehingga otomatis menyusut bila
%   lebih lebar dari teks; atur ukuran sumber lewat fig.width/fig.height (lihat index.qmd).
% - Tabel lebar: gunakan kolom 'X' (tabularx), kableExtra column_spec(width=...),
%   atau putar ke landscape. Jangan memperkecil huruf < 8 pt (pedoman).
% - Rumus panjang: penggal manual dgn aligned/split/multline (operator di AWAL baris lanjutan).

% --- Penomoran & gaya judul bab/subbab (Times, tebal) ---
%     Judul bab tampil "BAB I, BAB II, ..." (Romawi besar), tetapi penomoran
%     subbab tetap memakai angka Arab dengan pemisah titik (1.1, 1.2, 2.1, ...)
%     sesuai pedoman hlm. 34-35. Karena itu counter chapter dibiarkan Arab,
%     dan Romawi hanya ditampilkan pada JUDUL bab.
\usepackage{titlesec}
\renewcommand{\chaptername}{BAB}
% Judul bab/subbab diketik 1 spasi (pedoman): bungkus dengan \singlespacing.
\titleformat{\chapter}[display]
  {\normalfont\bfseries\centering\singlespacing}
  {\MakeUppercase{\chaptername}\ \Roman{chapter}}{0pt}
  {\MakeUppercase}
% Jarak judul bab ke teks/subbab di bawahnya (pedoman hlm. 34). Unit "1 spasi" =
% \normalbaselineskip (tinggi baris TUNGGAL, tidak terpengaruh \doublespacing teks
% isi), sehingga nilainya deterministik. Jarak bab->teks = 3 spasi (afterskip bab).
\titlespacing*{\chapter}{0pt}{0pt}{3\normalbaselineskip}
\titleformat{\section}{\normalfont\bfseries\singlespacing}{\thesection}{1em}{}
\titleformat{\subsection}{\normalfont\bfseries\singlespacing}{\thesubsection}{1em}{}
% Judul level ke-4 (####) tidak bernomor (number-depth:3) -> cukup TEBAL (pedoman
% hlm. 35), tanpa miring.
\titleformat{\subsubsection}{\normalfont\bfseries\singlespacing}{\thesubsubsection}{1em}{}
% Jarak judul DETERMINISTIK (unit = \normalbaselineskip = 1 spasi tunggal), sesuai
% pedoman hlm. 34-35:
%   - bab -> subbab    = afterskip bab (3) + beforeskip subbab (1) = 4 spasi
%   - bab/subbab -> teks = 3 spasi (afterskip subbab; afterskip bab juga 3)
%   - judul level-4 (pokok bahasan) -> 4 spasi di atasnya (beforeskip subsubsection)
% titlesec menjumlahkan afterskip bab + beforeskip subbab pada transisi bab->subbab.
\titlespacing*{\section}{0pt}{1\normalbaselineskip}{3\normalbaselineskip}
\titlespacing*{\subsection}{0pt}{1\normalbaselineskip}{3\normalbaselineskip}
\titlespacing*{\subsubsection}{0pt}{4\normalbaselineskip}{1\normalbaselineskip}

% Judul bab TAK bernomor (DAFTAR ISI, DAFTAR TABEL, DAFTAR GAMBAR, PRAKATA,
% ABSTRAK, DAFTAR PUSTAKA, dll.) dibuat seragam dengan judul bab bernomor:
% huruf tebal, rata tengah, ukuran normal (12pt), 1 spasi, tanpa "BAB".
\titleformat{name=\chapter,numberless}[display]
  {\normalfont\bfseries\centering\singlespacing}
  {}{0pt}
  {\MakeUppercase}
% Jarak judul bab tak-bernomor ke isi di bawahnya = 3 spasi (deterministik),
% seragam dengan bab bernomor.
\titlespacing*{name=\chapter,numberless}{0pt}{0pt}{3\normalbaselineskip}

% Judul daftar (DAFTAR ISI/TABEL/GAMBAR/LAMPIRAN) diseragamkan dengan judul
% bab — tebal, rata tengah, ukuran normal (12pt), jarak ~2 baris — melalui API
% resmi tocloft (\cftXtitlefont). Pengaturannya diletakkan setelah paket
% 'tocloft' dimuat (lihat blok Daftar Lampiran di bawah).

% --- Penomoran bab Romawi (BAB I) untuk JUDUL & DAFTAR ISI, dan Arab untuk
%     subbab (mis. 1.1, 1.1.1). TABEL, GAMBAR, dan PERSAMAAN memakai nomor urut
%     TUNGGAL berjalan di seluruh skripsi (Tabel 1, 2, 3, ...; Gambar 1, 2, ...;
%     persamaan (1), (2), ...) sesuai pedoman hlm. 5 & 8 (bukan per-bab 4.1).
%     Diletakkan pada \AtBeginDocument agar menimpa definisi bawaan Quarto.
\AtBeginDocument{%
  \renewcommand{\thechapter}{\Roman{chapter}}%
  \renewcommand{\thesection}{\arabic{chapter}.\arabic{section}}%
  \renewcommand{\thesubsection}{\arabic{chapter}.\arabic{section}.\arabic{subsection}}%
  \renewcommand{\thesubsubsection}{\arabic{chapter}.\arabic{section}.\arabic{subsection}.\arabic{subsubsection}}%
  % Lepas reset-per-bab -> nomor tabel/gambar/persamaan menjadi deret tunggal.
  \counterwithout{table}{chapter}%
  \counterwithout{figure}{chapter}%
  \counterwithout{equation}{chapter}%
}

% --- Nomor persamaan ber-kurung pada TAMPILAN dan pada RUJUKAN (\ref) ---
%     Dengan eq-prefix kosong di _quarto.yml, Quarto menulis rujukan persamaan
%     hanya sebagai "\ref{eq-...}" (tanpa kata "Persamaan"). \ref ke persamaan
%     mengembalikan \theequation; karena \theequation diberi kurung, rujukan
%     menjadi "(1)" saja, sehingga kata "persamaan" ditulis manual di naskah.
%     Kurung bawaan pada \tagform@ dihilangkan supaya tampilan persamaan tidak
%     menjadi ganda ((1)).
\makeatletter
\AtBeginDocument{%
  \renewcommand{\theequation}{(\arabic{equation})}%   nomor & rujukan: (n)
  \renewcommand{\tagform@}[1]{\maketag@@@{#1}}%        cegah kurung ganda di tampilan
}
\makeatother

% --- Penamaan caption: "Tabel" dan "Gambar" (bukan Table/Figure) ---
%     Mengatur teks pada float; awalan rujukan silang diatur di _quarto.yml.
\renewcommand{\tablename}{Tabel}
\renewcommand{\figurename}{Gambar}

% --- Caption tabel/gambar (Times); judul tabel di ATAS, judul gambar di BAWAH.
%     Judul tabel/gambar yang lebih dari satu baris diketik 1 spasi (pedoman). ---
\usepackage{caption}
\captionsetup{font={normalsize,stretch=1.0},labelsep=period,justification=centering}
\captionsetup[table]{position=top}
\captionsetup[figure]{position=bottom}

% --- Tata letak float (pedoman): tabel/gambar SELALU rata tengah dan berjarak
%     ~3 spasi dari teks di atas/bawahnya. Nilai \intextsep dll. di bawah
%     dihitung saat 1,5 spasi; karena teks isi kini 2 spasi, jaraknya hanya
%     mendekati 3 spasi — sesuaikan bila perlu setelah render. ---
\usepackage{float}
\makeatletter
\g@addto@macro\@floatboxreset\centering   % semua isi float otomatis rata tengah
\makeatother
\setlength{\intextsep}{3\normalbaselineskip}     % jarak float (here/H) dgn teks atas-bawah = 3 spasi
\setlength{\textfloatsep}{3\normalbaselineskip}   % jarak float atas/bawah halaman dgn teks = 3 spasi
\setlength{\floatsep}{2\normalbaselineskip}       % jarak antar-float
\setlength{\abovecaptionskip}{0.5\baselineskip}
\setlength{\belowcaptionskip}{0.5\baselineskip}

% --- Nomor halaman:
%     bagian awal  -> angka Romawi kecil, TENGAH bawah  (gaya 'plain' bawaan)
%     bagian isi   -> angka Arab, di KANAN bawah (pedoman hlm. 1), berlaku pada
%                     halaman ganjil maupun genap.
%     Karena \chapter memanggil \thispagestyle{plain}, kita sediakan
%     perintah \aktifkanNomorIsi yang mendefinisikan ulang 'plain' agar
%     nomor halaman bab maupun halaman biasa berada di kanan bawah.
\usepackage{fancyhdr}
\fancypagestyle{isistis}{%
  \fancyhf{}%
  \renewcommand{\headrulewidth}{0pt}%
  \fancyfoot[R]{\thepage}%  bagian isi: angka Arab di tepi kanan bawah (pedoman hlm. 1)
}
\newcommand{\aktifkanNomorIsi}{%
  \fancypagestyle{plain}{%
    \fancyhf{}%
    \renewcommand{\headrulewidth}{0pt}%
    \fancyfoot[R]{\thepage}%
  }%
  \pagestyle{isistis}%
}
\pagestyle{plain}

% --- Daftar Isi/Tabel/Gambar/Lampiran/Pustaka berhuruf kapital ---
%     Diletakkan pada \AtBeginDocument agar menimpa definisi yang
%     disuntikkan Quarto/babel belakangan (yang memakai huruf kapital
%     di awal saja, mis. "Daftar Tabel").
\AtBeginDocument{%
  \renewcommand{\contentsname}{DAFTAR ISI}%
  \renewcommand{\listfigurename}{DAFTAR GAMBAR}%
  \renewcommand{\listtablename}{DAFTAR TABEL}%
  \renewcommand{\bibname}{DAFTAR PUSTAKA}%
  \renewcommand{\refname}{DAFTAR PUSTAKA}%
}

% --- Daftar Lampiran (memakai mekanisme 'tocloft' agar serupa daftar lain) ---
\usepackage{tocloft}
\newcommand{\listlampiranname}{DAFTAR LAMPIRAN}
\newlistof{lampiran}{lmp}{\listlampiranname}
% \lampiran{Judul}: daftarkan satu lampiran ke Daftar Lampiran. Setiap lampiran
% BARU dimulai di halaman berikutnya (\clearpage); lampiran pertama mengikuti
% judul "LAMPIRAN" tanpa halaman kosong (syarat \ifnum: hanya pindah halaman
% jika sudah ada lampiran sebelumnya).
\newcommand{\lampiran}[1]{%
  \ifnum\value{lampiran}>0\relax\clearpage\fi
  \refstepcounter{lampiran}%
  \addcontentsline{lmp}{lampiran}{\protect\numberline{\thelampiran}#1}%
}

% Judul keempat daftar (ISI, TABEL, GAMBAR, LAMPIRAN) diseragamkan dengan
% judul bab: tebal, rata tengah, ukuran normal, jarak ~2 baris atas-bawah.
% Catatan: gunakan \hfil (bukan \hfill) agar judul benar-benar terpusat di
% tengah lebar teks. \hfill memiliki regangan tak hingga yang berinteraksi
% dengan glue internal tocloft sehingga judul meleset ke kanan pada halaman
% ganjil (tata letak bolak-balik).
\renewcommand{\cfttoctitlefont}{\hfil\normalfont\bfseries\MakeUppercase}
\renewcommand{\cftaftertoctitle}{\hfil\par\nobreak\vskip2\spasi\mbox{}\hfill{\normalfont Halaman}}
\renewcommand{\cftlottitlefont}{\hfil\normalfont\bfseries\MakeUppercase}
\renewcommand{\cftafterlottitle}{\hfil\par\nobreak\vskip2\spasi{\normalfont No.~Tabel\hfill Judul Tabel\hfill Halaman}}
\renewcommand{\cftloftitlefont}{\hfil\normalfont\bfseries\MakeUppercase}
\renewcommand{\cftafterloftitle}{\hfil\par\nobreak\vskip2\spasi{\normalfont No.~Gambar\hfill Judul Gambar\hfill Halaman}}
\renewcommand{\cftlmptitlefont}{\hfil\normalfont\bfseries\MakeUppercase}
\renewcommand{\cftafterlmptitle}{\hfil\par\nobreak\vskip2\spasi{\normalfont No.~Lampiran\hfill Judul Lampiran\hfill Halaman}}
\setlength{\cftbeforetoctitleskip}{2\baselineskip}\setlength{\cftaftertoctitleskip}{1.5\spasi}
\setlength{\cftbeforelottitleskip}{2\baselineskip}\setlength{\cftafterlottitleskip}{\spasi}
\setlength{\cftbeforeloftitleskip}{2\baselineskip}\setlength{\cftafterloftitleskip}{\spasi}
\setlength{\cftbeforelmptitleskip}{2\baselineskip}\setlength{\cftafterlmptitleskip}{1.5\baselineskip}

% --- Jarak entri Daftar Gambar & Daftar Tabel seragam (1 spasi) ---
%     report.cls menambah \addvspace{10pt} ke .lof/.lot di setiap pergantian
%     bab, sehingga entri gambar/tabel dari bab berbeda terlihat berjarak lebih
%     lebar daripada yang sebab. Hapus agar semua entri berjarak sama (pedoman:
%     daftar ditulis 1 spasi).
\makeatletter
\patchcmd{\@chapter}{\addtocontents{lof}{\protect\addvspace{10\p@}}}{}{}{}
\patchcmd{\@chapter}{\addtocontents{lot}{\protect\addvspace{10\p@}}}{}{}{}
\makeatother

% --- Judul elemen front-matter non-tocloft (mis. Daftar Singkatan & Lambang) ---
%     Meniru PERSIS gaya judul daftar tocloft: rata tengah, tebal, KAPITAL,
%     jarak ~2 baris atas-bawah, dan terdaftar di Daftar Isi sebagai entri bab.
%     Memakai \clearpage (bukan \cleardoublepage) agar tidak menyisipkan halaman
%     kosong, konsisten dengan daftar lain. #1 = judul tampil, #2 = entri Daftar Isi.
\newcommand{\judulfrontmatter}[2]{%
  \clearpage
  \phantomsection
  \addcontentsline{toc}{chapter}{#2}%
  \vspace*{2\baselineskip}%
  {\centering\normalfont\bfseries\MakeUppercase{#1}\par}%
  \vspace{2\baselineskip}%
}

% --- Daftar Isi: entri bab tampil "BAB I  PENDAHULUAN ........ hlm." ---
%     Awalan "BAB" ditambahkan pada entri bab bernomor; bab tak bernomor
%     (Prakata, Abstrak, Daftar Pustaka, Lampiran, Riwayat Hidup) tetap tanpa
%     awalan. Diberi pemandu titik (dotted leader) dan bobot normal.
\renewcommand{\cftchapfont}{\normalfont}
\renewcommand{\cftchappagefont}{\normalfont}
\renewcommand{\cftchapleader}{\normalfont\cftdotfill{\cftdotsep}}
\renewcommand{\cftchappresnum}{BAB\space}
\renewcommand{\cftchapaftersnum}{\hspace{0.5em}}
\setlength{\cftchapnumwidth}{5.2em}   % cukup untuk "BAB VIII"

% --- Halaman genap yang tersisip otomatis (akibat openright/twoside)
%     diberi keterangan "Halaman ini sengaja dikosongkan" di tengah,
%     sesuai konvensi pedoman. Mendefinisikan ulang \cleardoublepage. ---
\renewcommand{\cleardoublepage}{%
  \clearpage
  \ifodd\value{page}\else
    \thispagestyle{empty}%
    \vspace*{\fill}
    \begin{center}\itshape Halaman ini sengaja dikosongkan\end{center}
    \vspace*{\fill}
    \clearpage
  \fi
}

% --- Nonaktifkan halaman judul bawaan Quarto (\maketitle) karena halaman
%     muka STIS sudah disisipkan melalui tex/00_frontmatter.tex ---
\AtBeginDocument{\renewcommand{\maketitle}{}}

% ============================================================
%  TikZ — untuk diagram alir (flowchart), use case diagram, dan
%  gambar vektor lain. Caption gambar otomatis "Gambar X.Y" (lihat
%  pengaturan \figurename & \thefigure di atas).
% ============================================================
\usepackage{tikz}
\usetikzlibrary{shapes.geometric,shapes.symbols,shapes.misc,arrows.meta,positioning,calc,fit,backgrounds,chains}

% --- Gaya simpul diagram alir (flowchart) ---
\tikzset{
  % Terminator (Mulai/Selesai) — kapsul
  fc-mulai/.style   = {rounded rectangle, draw, thick, minimum width=2.6cm,
                       minimum height=0.9cm, align=center, font=\normalsize},
  % Proses — persegi panjang
  fc-proses/.style  = {rectangle, draw, thick, minimum width=3cm,
                       minimum height=1cm, align=center, font=\normalsize},
  % Input/Output — jajar genjang
  fc-io/.style      = {trapezium, trapezium left angle=70, trapezium right angle=110,
                       draw, thick, minimum width=2.6cm, minimum height=0.9cm,
                       align=center, font=\normalsize},
  % Keputusan — belah ketupat
  fc-putusan/.style = {diamond, draw, thick, aspect=2, inner sep=2pt,
                       minimum width=2.8cm, align=center, font=\normalsize},
  % Anak panah alur
  fc-panah/.style   = {-{Stealth[length=2.5mm]}, thick},
}

% --- Gaya elemen use case diagram ---
\tikzset{
  % Aktor (stick figure) sederhana
  uc-aktor/.pic = {
    \draw[thick] (0,0) circle (0.18);                  % kepala
    \draw[thick] (0,-0.18) -- (0,-0.75);               % badan
    \draw[thick] (-0.35,-0.40) -- (0.35,-0.40);        % tangan
    \draw[thick] (0,-0.75) -- (-0.28,-1.15);           % kaki kiri
    \draw[thick] (0,-0.75) -- (0.28,-1.15);            % kaki kanan
  },
  % Use case — elips
  uc-case/.style = {ellipse, draw, thick, align=center,
                    minimum width=2.8cm, minimum height=1cm, font=\small},
  % Garis asosiasi aktor–use case
  uc-asosiasi/.style = {thick},
  % Relasi <<include>> / <<extend>> (panah putus-putus)
  uc-relasi/.style   = {-{Stealth[length=2.5mm]}, thick, dashed, font=\footnotesize},
  % Kotak batas sistem (system boundary)
  uc-sistem/.style   = {rectangle, draw, thick, rounded corners=2pt,
                        inner sep=0.6cm},
}

% ============================================================
%  Algoritma (algorithm2e) — penamaan & kata kunci bahasa Indonesia,
%  penomoran deret tunggal (mis. Algoritme 1, 2, ...).
% ============================================================
\usepackage[ruled,vlined,linesnumbered]{algorithm2e}
\SetAlgorithmName{Algoritme}{Algoritme}{Daftar Algoritme}
% Kata kunci dalam bahasa Indonesia
\SetKwInput{KwIn}{Masukan}
\SetKwInput{KwOut}{Keluaran}
\SetKwInput{KwData}{Data}
\SetKwInput{KwResult}{Hasil}
\SetKwFor{For}{untuk}{lakukan}{akhir untuk}
\SetKwFor{ForEach}{untuk setiap}{lakukan}{akhir untuk}
\SetKwFor{While}{selama}{lakukan}{akhir selama}
\SetKwIF{If}{ElseIf}{Else}{jika}{maka}{jika lain}{lainnya}{akhir jika}
\SetKwBlock{Begin}{mulai}{akhir}
\SetKw{KwTo}{sampai}
\SetKw{KwRet}{kembalikan}
\SetKw{Return}{kembalikan}
% Penomoran algoritme: deret tunggal (Algoritme 1, 2, ...) agar konsisten
% dengan tabel/gambar/persamaan. (Pedoman tidak mengatur algoritme.)
\renewcommand{\thealgocf}{\arabic{algocf}}

% ============================================================
%  Daftar Singkatan & Lambang berbasis paket glossaries (opsional).
%  Diisi otomatis oleh buat_skripsi() sesuai argumen `daftar_singkatan`:
%   - mode "glossaries": memuat paket + entri (otomatis muncul bila dipakai
%     lewat \gls{...}); memakai \makenoidxglossaries (murni LaTeX, tanpa
%     makeglossaries eksternal).
%   - mode lain: hanya menyediakan \gls dkk. sebagai jaring pengaman agar
%     dokumen tetap kompilasi bila ada sisa \gls{...}.
% ============================================================
{{GLOSSARIES_PREAMBLE}}

% ============================================================
%  Kustomisasi mahasiswa (TIDAK ditimpa saat pembaruan paket).
%  Tulis paket / perintah LaTeX tambahan Anda di berkas terpisah:
%  tex/preamble-tambahan.tex
% ============================================================
\IfFileExists{tex/preamble-tambahan.tex}{\input{tex/preamble-tambahan.tex}}{}
